% 14M control: Run048 observations/001-causal-hz-latency.md; 07_reduce.py.
% Evidence and editorial scope: reviews/2026-09-21-results-arc/README.md (approved abstract structure).
% Analysis024 observations 015, 028, 030, 034; Analysis021 pressure-scope O016.
% 70M adaptation: Analysis028 observations/001-matched-70m-grid.md; Run042 controls and Run045 grid.
\begin{abstract}
% Training-time activation pressure ($P$) induce smaller magnitudes, while thresholding ($T$) maps small values to exact zeros. Where should these interventions be placed to reduce inference latency at an acceptable quality cost? We study broad thresholding recipes ($T_4,T_7$) with local or global pressure under matched Pythia-14M pretraining conditions, coupling sparsity measurements with specialized GPU execution. \textbf{Higher model-wide sparsity does not guarantee lower latency}, and broader interventions can incur greater quality costs.

% \textbf{Sparsity spillover and execution diagnostics converge on the same sites.} Pressure at the FFN hidden activation ($h$) also increases sparsity at the unpressured attention-output activation ($z$). Controlled kernel ablations identify these sites as the profitable execution paths, with conditional savings of approximately 150 and 29\,$\mu$s, respectively, on $T_7/P_{\mathrm{all}}$ at $\kappa=0.5$. This motivates a narrower intervention, $T_2/P_h$, that pressures $h$ and thresholds only $h,z$. \textbf{Targeted placement preserves quality while exposing useful sparsity}: at $\kappa=0.1$, validation loss is 5.151 versus Base's 5.209, with latency 0.573 versus 0.655\,ms for native Base (RTX\,5090, batch one, 2,048 tokens).

% At 70M, $T_2/P_h$ retains the lowest loss among the tested threshold topologies, but its favorable quality--latency balance does not transfer directly. The kernel requires further optimization; a same-checkpoint control on $T_7/P_h$ at $\kappa=0.5$ then shows a 17.4\% latency reduction from specialized $h,z$ execution relative to native dense replacements. This conditional benefit includes fusion, layout and inspection costs alongside sparsity exploitation. Training interventions and deployment kernels must therefore be co-designed so that sparsity is induced where it is structurally exploitable at inference time.

% Training-time activation pressure ($P$) induces smaller activation magnitudes, while thresholding ($T$) maps small values to exact zeros. Where should these interventions be placed to reduce inference latency at an acceptable validation-loss cost? We study broad thresholding recipes ($T_4,T_7$) with local or global pressure under matched Pythia-14M pretraining conditions, coupling activation-sparsity measurements with specialized GPU execution. We find that higher model-wide sparsity does not necessarily yield lower latency, and that broader interventions can incur substantially larger loss penalties. Sparsity and execution diagnostics instead converge on two sites. Pressure applied only at the FFN hidden activation ($h$) also increases sparsity at the unpressured attention-output activation ($z$), a nonlocal response we call sparsity spillover. Conditional kernel ablations identify $h$ and $z$ as the profitable skipping paths. This motivates a narrower intervention, $T_2/P_h$, which pressures $h$ and thresholds only $h,z$. At moderate threshold ($\kappa=0.1$), it reaches validation loss 5.151 versus 5.209 for the base model. With weights and thresholding fixed, exploiting its $h,z$ zeros reduces latency by $12.9\%$ relative to the same implementation with both skipping paths disabled. At 70M, $T_2/P_h$ retains the lowest validation loss among the tested threshold topologies, but the 14M optimized kernel does not transfer directly. After further kernel optimization, $T_7/P_h$ at $\kappa=0.5$ shows a 17.4\% latency reduction from specialized $h,z$ execution relative to native dense replacements. These results show that activation sparsification is a training--execution co-design problem: sparsity must be induced where its structure can be profitably exploited at inference time.

% Scale extension: analyses/038-2026-09-24-base-t2-t7-scale-frontier/observations/003-manuscript-integration.md; 04_manuscript_evidence.py.
Activation sparsity can accelerate transformer inference when a specialized kernel can skip computation associated with zero activations. We ask how broadly and aggressively train-time sparsification should be applied to obtain such exploitable sparsity without excessive quality loss. Using matched Pythia-family pretraining runs, we systematically vary the placement and intensity of activation pressure and fixed-threshold nonlinearities, and connect the resulting sparsity to specialized-kernel execution. We find that broader interventions can substantially increase model-wide sparsity without reducing latency: at 14M parameters, more than doubling model-wide sparsity from 12.71\% to 27.48\% slightly increases latency. At 14M, conditional kernel ablations identify the FFN hidden activation \(h\) and attention-output projection input \(z\) as the most useful sparsification sites. Pressure applied only at \(h\) also increases sparsity at \(z\), despite \(z\) receiving no direct pressure, while broader pressure produces stronger optimization conflict at less responsive sites. These observations motivate thresholding only \(h,z\) while applying pressure only at \(h\). At 14M, this intervention achieves lower validation loss than the dense control while reducing latency from 0.655 to 0.573 ms. Comparisons at 31M and 70M retain the lower loss of targeted versus seven-site thresholding at matched thresholds, but not the validation-loss improvement over the dense control. Our results suggest that sparsification is best treated as a training--inference co-design problem: rather than maximizing sparsity globally, training should induce specific exploitable structure for the inference system.

\end{abstract}
